\documentclass{article}

\textwidth=6.5in
\textheight=8.5in
\voffset=-54pt
\oddsidemargin=0in
\evensidemargin=0in
\setlength{\parskip}{8pt}
\setlength{\parindent}{0pt}

\usepackage{amsmath, amssymb}
\usepackage{ifthen}

\usepackage[inline]{enumitem}
\setlist{topsep=0pt, parsep=8pt}

\usepackage[rgb,table]{xcolor}\convertcolorsUfalse
\usepackage{graphicx}

% change hyperref options
\PassOptionsToPackage{hyphens}{url}
\usepackage[bookmarks,colorlinks,linkcolor=black,citecolor=blue,pdfstartview=FitH,urlcolor=blue]{hyperref}
\hypersetup{pdfpagemode=UseNone}
\usepackage{bookmark}

\usepackage{graphicx}
\usepackage{multicol}
\usepackage{framed}
\usepackage{array}

\usepackage{icomma}

\usepackage{pifont}

\usepackage{soul}
\usepackage[normalem]{ulem}

\usepackage{tikz}
\usetikzlibrary{
  matrix,%
  % enables the snake paths
  decorations.pathmorphing,%
  % enables bigger arrows
  decorations.markings,%
  decorations.pathreplacing,%
  decorations.text,%
  calc,%
  patterns,%
  shapes.geometric,%
  arrows,
  decorations.shapes,
  positioning,
  plotmarks,
  intersections
}
\tikzset{>=latex} % sets default arrow type in tikz
\tikzset{font=\normalsize}
\tikzset{mark size=1.5pt, mark options=thin}
\tikzset{pin distance=4pt,
  every pin edge/.style={<-, thin, shorten <= -2pt}}

\interfootnotelinepenalty=10000
\renewcommand{\thefootnote}{(\arabic{footnote})}

\usepackage{multirow, bigdelim}

% change to \usepackage[sols]{handout} to see solutions
\usepackage{handout}

\newcommand\iscurrentenumi[1]{%
  \ifthenelse{\equal{#1}{\theenumi}}{}{#1}}
\setenumerate[1]{ref=\#\arabic*}
\setenumerate[2]{ref=\protect\iscurrentenumi{\theenumi}(\alph*)}
\newcommand\enumiiprefix[2]{%
  \ifthenelse{{\equal{#1}{\theenumi}}\AND{\equal{#2}{(\alph{enumii})}}}{}{}%
  \ifthenelse{{\equal{#1}{\theenumi}}\AND{\NOT{\equal{#2}{(\alph{enumii})}}}}{#2}{}%
  \ifthenelse{\equal{#1}{\theenumi}}{}{#1#2}%
}
\setenumerate[3]{ref=\protect\enumiiprefix{\theenumi}{(\alph{enumii})}\roman*}

\definecolor{purple}{rgb}{0.5,0,1.0}

\usepackage{amsthm}

% make URLs print in the same font
\usepackage{url}
\urlstyle{same}

\usepackage{pifont}
\def\exend{\hfill\ding{118}}

%\makeatletter\@addtoreset{equation}{section}\makeatother
%\renewcommand{\theequation}{\arabic{section}.\arabic{equation}}

\newtheoremstyle{theorem}% name
{8pt}%      Space above
{0pt}%      Space below
{\itshape}%         Body font
{}%         Indent amount (empty = no indent, \parindent = para indent)
{\bfseries}% Thm head font
{.}%        Punctuation after thm head
{.5em}%     Space after thm head: " " = normal interword space;
                                %       \newline = linebreak
{}%         Thm head spec (can be left empty, meaning `normal')

\theoremstyle{theorem}

\let\Theorem\undefined
\newtheorem{Theorem}[equation]{Theorem}
\let\Definition\undefined
\newtheorem{Definition}[equation]{Definition}
\let\Summary\undefined
\newtheorem{Summary}[equation]{Summary}
\let\Conjecture\undefined
\newtheorem{Conjecture}[equation]{Conjecture}
\let\Fact\undefined
\newtheorem{Fact}[equation]{Fact}

\renewcommand{\thesubsection}{\arabic{subsection}}
\makeatletter
\def\@seccntformat#1{\@ifundefined{#1@cntformat}%
   {\csname the#1\endcsname\quad}%       default
   {\csname #1@cntformat\endcsname}}%    enable individual control
\newcommand\section@cntformat{}
\makeatother

\usepackage{fancyhdr}
\lhead{\textcolor{white}{This content is protected and may not be shared, uploaded, or distributed.}}
\rhead{}
\rfoot{\vspace*{\fill}\footnotesize\textcopyright{} \emph{Harvard Math Ma}}
\renewcommand{\headrulewidth}{0pt}
\pagestyle{fancy}
\begin{document}

\begin{center}
  \Large\textsc{Final Exam Review Questions}  
\end{center}

We encourage you to make a study guide for yourself as you prepare for the final. Here are some questions to help guide you as you do this.

\textbf{Disclaimer:} This is \uline{not} meant to be a complete list of things we've studied. 

\begin{enumerate}
\item What do the graphs of $y = mx + b$, $y = x^2$, $y = x^3$, $y = \frac{1}{x}$, $y = \frac{1}{x^2}$, $y = |x|$, $y = \sqrt{x}$, $y = b^x$, and $y = \log_b x$ look like? (How do the graphs depend on $m$ and $b$?) These graphs are ones you want to just know.

\item How do you recognize graph transformations? For example, what kind of formula would give you a vertical stretch of $y = 2^x$? A horizontal stretch? A vertical shift? A horizontal shift?

\item We've looked at several characteristics of functions, like whether they're even, odd, or neither; where they're increasing or decreasing; where they're concave up or concave down; and whether they have horizontal and vertical asymptotes. What do these characteristics look like graphically? How do we find them algebraically? 

\item What's the difference between a function being \emph{continuous} and a function being \emph{differentiable}? What are the different types of discontinuities? How do you find the discontinuities (and which type each one is) of a function given by an algebraic formula? 

\item What are different methods for calculating limits? When does each apply? How do you visualize limits? 

\item What's the difference between an average rate of change and an instantaneous rate of change? How do we visualize these on a graph?

\item What's the limit definition of the derivative? What does each piece of it mean? 

\item What's the idea of secant line approximation? Tangent line approximation? How are these similar, and how are they different? What kind of information do you need to be able to use each method? What are examples of questions you can answer using these methods? 

\item What are the important relationships between $f$, $f'$, and $f''$?

\item We've learned several differentiation rules (the Sum Rule, Constant Multiple Rule, Power Rule, Quotient Rule, and Chain Rule) as well as derivatives of several ``basic'' functions ($x^n$, $e^x$, $e^{kx}$, $b^x$, $\ln x$, $\log_b x$). What do all of these rules say? How do you recognize when to use each?

\item What's the difference between a local minimum and a global minimum? How do you find local mins? How about global?

\item What's an inverse function? How do you tell whether a function has an inverse function? What are examples of functions and their inverses? 

\item We've studied several types of functions this semester: linear, exponential, logarithmic, polynomial. What are key properties of each?

\item What types of situations might you model with a linear function? An exponential function? If you have two data points, how would you write a linear function through them? How about an exponential function? 

\item What's a logarithm?

\item There are 3 exponent rules and 3 logarithm rules to remember; what are they? 

\item What does the degree of a polynomial tell you about its graph?

\end{enumerate}

\end{document}
